% !TeX root = ../thesis.tex
\chapter{Introduction}

Digital tools have become a standard part of education. From primary school to university, students interact with adaptive practice platforms, intelligent tutoring systems, learning management environments, and a growing number of AI-driven applications \cite{dancsa_digital_2023}. In primary school arithmetic alone, platforms such as Automatus \cite{automatus} and Rekentuin \cite{rekentuin} are used daily by hundreds of thousands of children across the Netherlands and Flanders. These tools have shown that \textbf{adaptive practice, gamification, and child-friendly design} can sustain engagement and support the development of arithmetic fluency.

In parallel, educational psychology researchers rely on a different category of digital tools to investigate \emph{how} children learn arithmetic. Platforms such as OpenSesame \cite{mathot_opensesame_2012} and PsychoPy \cite{peirce2019psychopy2} are designed to support experimental research, offering \textbf{precise control over stimulus presentation, complete data logging, and reproducible study configurations.} These tools have made it possible to study aspects of arithmetic development such as fact retrieval, strategy use, and metacognitive monitoring \cite{bellon_more_2019, rinne_knowing_2014}.

The problem is that \textbf{these two categories of tools have evolved in isolation.} Classroom tools are built to engage children but are not built for research. Research tools are built for scientific control but are not built for classrooms. A researcher who wants to study arithmetic learning in a real classroom, often across a training period of several sessions rather than a single visit, is therefore forced to choose between two unsatisfactory options. The first is to use a research tool: in principle these tools can be extended with stories, rewards, and adaptivity, but doing so requires substantial programming, and in practice children are left with a sparse interface. The second is to use a classroom tool and accept that the experimental conditions cannot be controlled and the relevant data cannot be collected. Neither choice is good, and the compromise that results is rarely satisfactory.

This compromise is not just inconvenient. It directly affects the quality of the research that comes out of these studies. Arithmetic research with primary school children typically takes place over several sessions, spread across days or weeks, and depends on children staying motivated and focused throughout. When the tool used is unfamiliar and visually bare, three things tend to go wrong. First, children waste time and mental energy figuring out how the tool works rather than focusing on the actual exercises. Second, sitting in front of a screen with no rewards, no story, and no sense of progress is exhausting, and children's performance drops as the sessions go on \cite{sievertsen2016cognitive, hopstaken2015role}. Third, the artificial setting of a stripped-down experiment looks nothing like the classroom the research is ultimately trying to understand, which limits how much the findings actually transfer back to real-world learning \cite{bronfenbrenner1977toward}. \textbf{A better tool, in short, means better research.}

This thesis introduces \textbf{RekenRangers (MathRangers)}, a web-based arithmetic learning platform designed to bridge the gap between classroom exercise tools and controlled research software. RekenRangers offers a gamified, child-friendly environment built around a safari theme, while also giving researchers fine-grained control over experimental conditions, complete data logging, and the ability to configure studies \textbf{without writing a single line of code.} The goal is a single tool that children enjoy using and that researchers can rely on for their scientific work.

RekenRangers was developed in close collaboration with the co-supervisor of this thesis, Elien Bellon, and with two master's students in pedagogical sciences at KU Leuven, Eline T'joen and Kiara Willems. Their thesis project required exactly the kind of tool this work set out to build: a platform that could run a multi-condition arithmetic study in real classrooms while keeping children motivated across multiple training sessions. This collaboration shaped many of the design decisions and made it possible to evaluate the tool in practice, not just in theory.

To test whether RekenRangers actually delivers on its promise, the platform was deployed in a real classroom study with 107 third-grade children across seven primary schools in Flanders. Children worked with the tool over four training days, and afterwards filled in a short questionnaire about their experience. The two researchers who ran the study also provided written feedback on what it was like to use RekenRangers to set up, run, and monitor a real experiment. Together, the data, the questionnaire results, and the researcher feedback form the basis for the evaluation presented later in this thesis.

The results of this evaluation are encouraging, but with clear caveats. RekenRangers proved capable of supporting a real multi-session classroom experiment from start to finish: 21,522 exercise records were collected across four experimental conditions, the data was structured and exportable without preprocessing, and children remained visibly engaged with the tool throughout. The behavioural data shows that 98.5\% of earned coins were spent voluntarily on puzzles and that 99\% of participants completed at least one puzzle, both of which suggest that the gamification system was not just decoration but a genuine source of motivation. The questionnaire results support this picture, with children rating the tool positively on engagement and user friendliness.

At the same time, not everything worked as designed. The adaptive difficulty system differentiated between children of different ability levels but did not converge on its intended target accuracy, indicating that the algorithm needs further calibration. Two related issues surfaced in the difficulty data: a ceiling effect, with several children reaching the maximum rating because the exercise pool did not extend far enough at the top end, and a plateau, with a sizeable group of children getting stuck around the boundary between two difficulty zones because the step in difficulty between those zones was too large. A submission bug led to a small number of duplicate records that had to be cleaned during analysis. And while the configuration of an experiment can be done without writing code, the interface for doing so still has room for improvement, particularly in helping researchers see and verify the experiment they are setting up.

Some properties of the tool are present but underdeveloped. Personalisation is functional but lacks the visual elements (avatars, personal dashboards) that would make the experience feel truly personal. The feedback system is minimal, providing only correctness information. And the flexibility of the tool, while sufficient for arithmetic studies that fit its built-in templates, cannot match the open-ended configurability of a programming environment. These are the natural directions for future work.

Taken together, the work shows that bridging the gap between classroom exercise tools and controlled research software is achievable in principle, while making the remaining trade-offs and open problems explicit. RekenRangers is not a finished product, but it is \textbf{a working demonstration that the two worlds can be brought closer than they currently are.}

The remainder of this thesis is structured as follows. Chapter~\ref{chp:background} reviews arithmetic research and existing research tools. Chapter~\ref{chp:exercise-tools} does the same for classroom exercise tools. Chapter~\ref{chp:problem-statement} formalises the gap and defines the research questions. Chapter~\ref{chp:rekenrangers} introduces RekenRangers and its design. Chapter~\ref{chp:experimental-setup} describes the classroom deployment. Chapter~\ref{chp:results} presents the evaluation results. Chapter~\ref{chp:discussion} interprets these results and discusses the limitations. Chapter~\ref{chp:conclusions} answers the research questions and outlines future work.

\vspace{5em}
\noindent
\begin{tcolorbox}[
    colback=gray!5!white,
    colframe=gray!50!black,
    title=Note on the use of generative AI
]
Generative AI tools were used in the preparation of this thesis. Their use was limited to assisting with writing, specifically the reformulation and stylistic refinement of text that had already been drafted by the author, and to support the search for relevant scientific literature. Any sources surfaced in this way were always retrieved, read, and critically assessed by the author before being cited, and no claim or reference was included on the basis of generative AI output alone. All scientific content, design decisions, analyses, and conclusions are the author's own. Generative AI was also used during the development of RekenRangers to assist with implementation tasks such as generating boilerplate code and debugging, as well as to create the visual assets used in the safari theme.
\end{tcolorbox}

